\documentclass{article}
\usepackage[utf8]{inputenc}
\usepackage{amsmath}
\usepackage{geometry}
\geometry{margin=2cm}
\begin{document}

\section*{Análise da Questão 160 - ENEM 2024 - Matemática}

\subsection*{Enunciado}

Uma piscina tem capacidade de 2 500 000 litros. Seu sistema de abastecimento foi regulado para ter uma vazão constante de 6 000 litros de água por minuto. O mesmo sistema foi instalado em uma segunda piscina, com capacidade de 2 750 000 litros, e regulado para ter uma vazão, também constante, capaz de enchê-la em um tempo 20\% maior que o gasto para encher a primeira piscina.

A vazão do sistema de abastecimento da segunda piscina, em litro por minuto, é:

\subsection*{Resolução Técnica}

A vazão está relacionada ao volume e ao tempo pelo seguinte:

\[
vazão = \frac{volume}{tempo}.
\]

Para a primeira piscina, temos:

\[
  V_1 = 2\;500\;000 \text{ litros}, \quad Q_1 = 6\;000 \text{ litros/minuto}.
\]

Calculamos o tempo para encher a primeira piscina:

\[
  t_1 = \frac{V_1}{Q_1} = \frac{2\;500\;000}{6\;000} = 416,67 \text{ minutos}.
\]

Como a segunda piscina é enchida em um tempo 20\% maior:

\[
  t_2 = t_1 \times 1,20 = 416,67 \times 1,20 = 500 \text{ minutos}.
\]

Sua capacidade é:

\[
  V_2 = 2\;750\;000 \text{ litros}.
\]

Logo, a vazão do sistema da segunda piscina é:

\[
  Q_2 = \frac{V_2}{t_2} = \frac{2\;750\;000}{500} = 5\;500 \text{ litros por minuto}.
\]

\textbf{Resposta:} 5 500 litros por minuto.

\subsection*{Explicação Didática}

A vazão indica quantos litros de água fluem por minuto. Conforme o tempo para encher aumenta em 20\%, a vazão diminui, pois o fornecimento de água é mais lento, apesar do volume ser maior. O problema ensina a aplicar a relação de grandezas inversamente proporcionais e destacar que o aumento no tempo implica redução na vazão.

\subsection*{Erros Comuns (Distratores)}

\begin{itemize}
  \item Supor que o aumento de 20\% no tempo implica aumento na vazão (como nos valores 8250 ou 7920). Esta é a concepção errada da proporcionalidade.
  \item Aplicar o percentual incorreto ou diretamente na vazão, errado por não considerar a relação inversa.
  \item Arredondar mal ou calcular errado diferenças percentuais (como 6545 ou 5280).
\end{itemize}

\subsection*{Perfil da Questão}

A questão exige raciocínio proporcional, habilidade com unidades e interpretação do texto para identificar que o aumento percentual ocorre no tempo e não na vazão. Tem dificuldade média e não requer elementos visuais ou modelagens complexas.

\bigskip

\textbf{Conclusão:} A análise é consistente e aprovada, trazendo todas as informações necessárias para a compreensão completa do problema.

\end{document}